\documentclass[../../../../uem_mat_mei_q.tex]{subfiles}

\begin{document}
    平面上に異なる4点O，A，B，Cがある．点Cは線分ABを$2:3$に内分する点である．%
    $\overrightarrow{\mathrm{OA}}=\vecb{a}$，%
    $\overrightarrow{\mathrm{OB}}=\vecb{b}$，%
    $\overrightarrow{\mathrm{OC}}=\vecb{c}$とする．%
    $\vecb{a}$，$\vecb{b}$，$\vecb{c}$は %
    $\bigl|\vecb{a}\bigr|=1$，$\bigl|\vecb{b}\bigr|=1$，%
    $\bigl|\vecb{c}\bigr|=\myfraca{1}{2}$を満たす．%
    
    \begin{myenub}
        \item $\vecb{c}$は

            \myspaceb%
            $\displaystyle%
                \vecb{c}=\myfracc{\myboxb{ア}}{\myboxb{イ}}%
                \vecb{a}+\myfracc{\myboxb{ウ}}{\myboxb{エ}}%
                \vecb{b}
            $\\
            と表せる．
        \item 実数$x$についての2次関数

            \myspaceb%
            $\displaystyle%
                f(x)=\bigl|\vecb{a}+x\vecb{b}\bigr|^2%
            $\\
            を考える．$f(x)$は
            
            \myspaceb%
            $\displaystyle%
                x=\myfracc{\myboxb{オ}}{\myboxb{カキ}}%
            $\\
            のとき最小となる．
    \end{myenub}
\end{document}
